\rtmtight
\begin{tblr}{width=\textwidth,colspec={|Q[c,m,2.1cm]|X[l,m]|},hlines,vlines,hspan=minimal,rowsep=1.5pt,colsep=3pt,row{1}={bg=headgray,font=\bfseries}}
\SetCell{c,m}\hfil\logo\hfil & \textbf{POLITEKNIK NEGERI MALANG}\par\textbf{JURUSAN TEKNOLOGI INFORMASI}\par\textbf{PROGRAM STUDI: D4 TEKNIK INFORMATIKA} \\
\SetCell[c=2]{l} \textbf{RENCANA TUGAS MAHASISWA (RTM) DAN RUBRIK PENILAIAN} & \\
Nama Mata Kuliah & Jaringan Komputer \\
Kode Mata Kuliah & RTI254005 \quad \textbf{sks / Jam perminggu:} 2 SKS/4 jam \quad \textbf{Semester:} 4 \\
Dosen Pengampu & Tim Pengajar Program Studi D4 TI \\
\end{tblr}
\assessmentblock{Tugas Analisis Protokol Jaringan}{Case Method dan praktik konfigurasi jaringan}{SCPMK0504-03201, SCPMK0901-03202, SCPMK1002-03203}{Mahasiswa menyelesaikan tugas menganalisis arsitektur jaringan, mengkonfigurasi protokol, dan menganalisis aspek keamanan sesuai Sub-CPMK dan konteks mata kuliah.}{Menganalisis topologi/protokol, mengkonfigurasi simulator jaringan, mengidentifikasi masalah keamanan, dan mendokumentasikan hasil.}{Laporan analisis, konfigurasi jaringan, dan/atau bahan presentasi.}{Ketepatan analisis arsitektur dan protokol & 10 & Arsitektur jaringan dan protokol pada setiap lapisan dianalisis dengan benar, disertai justifikasi teknis yang kuat. & Analisis sebagian besar benar, tetapi justifikasi teknis atau detail lapisan masih kurang lengkap. & Analisis tidak tepat atau tidak menunjukkan pemahaman arsitektur jaringan dan protokol. \\ \\
Kualitas konfigurasi jaringan & 10 & Konfigurasi jaringan pada simulator berjalan dengan benar, semua protokol berfungsi, dan topologi sesuai spesifikasi. & Konfigurasi sebagian besar berjalan, tetapi beberapa protokol atau parameter belum dikonfigurasi dengan benar. & Konfigurasi tidak berjalan atau banyak parameter yang salah. \\ \\
Analisis keamanan dan dokumentasi & 10 & Ancaman keamanan diidentifikasi secara komprehensif, rekomendasi mitigasi tepat, dan dokumentasi lengkap. & Ancaman keamanan cukup diidentifikasi, tetapi rekomendasi atau dokumentasi masih terbatas. & Analisis keamanan tidak ada atau tidak relevan dengan topologi jaringan yang dianalisis. \\}

\assessmentblock{Kuis 1 Model OSI dan TCP/IP}{Kuis}{SCPMK0901-03202}{Kuis tertulis untuk mengukur pemahaman model OSI, TCP/IP suite, dan komponen jaringan.}{Menjawab soal pilihan ganda dan/atau uraian tentang model OSI dan TCP/IP.}{Jawaban tertulis kuis.}{Pemahaman model OSI dan fungsi lapisan & 3 & Ketujuh lapisan OSI, fungsinya, dan protokol yang berjalan di setiap lapisan dijelaskan dengan benar. & Sebagian besar lapisan dipahami, tetapi beberapa fungsi atau protokol masih kurang tepat. & Pemahaman model OSI tidak memadai atau banyak fungsi lapisan yang salah. \\ \\
Ketepatan TCP/IP dan komponen jaringan & 2 & TCP/IP suite, perbedaan dengan OSI, dan komponen jaringan (router, switch, hub) dijelaskan dengan benar dan lengkap. & TCP/IP cukup dipahami, tetapi perbandingan dengan OSI atau peran komponen masih kurang tepat. & TCP/IP atau komponen jaringan tidak dipahami dengan benar. \\}

\assessmentblock{UTS Protokol dan Lapisan Jaringan}{UTS berbasis analisis}{SCPMK0901-03202, SCPMK0504-03201}{Evaluasi tengah semester mengukur kemampuan menganalisis protokol jaringan dan lapisan komunikasi.}{Menganalisis skenario jaringan, menjawab soal analisis protokol, atau mengkonfigurasi jaringan sederhana sesuai soal.}{Jawaban analisis tertulis atau konfigurasi jaringan UTS.}{Analisis komponen dan protokol jaringan & 7 & Komponen jaringan dan protokol komunikasi dianalisis dengan benar sesuai skenario yang diberikan. & Analisis cukup benar, tetapi beberapa aspek protokol atau komponen belum dianalisis dengan tepat. & Analisis tidak tepat atau tidak menunjukkan pemahaman protokol jaringan. \\ \\
Konfigurasi dan penerapan protokol & 6 & Protokol dikonfigurasi dengan benar pada lapisan yang tepat dan menghasilkan komunikasi yang sesuai spesifikasi. & Konfigurasi sebagian besar benar, tetapi beberapa parameter atau lapisan masih belum dikonfigurasi optimal. & Konfigurasi tidak menghasilkan komunikasi yang benar atau tidak sesuai spesifikasi. \\}

\assessmentblock{Kuis 2 Keamanan Jaringan Dasar}{Kuis}{SCPMK1002-03203}{Kuis untuk mengukur pemahaman ancaman keamanan jaringan dan mekanisme perlindungan dasar.}{Menjawab soal tentang jenis ancaman, serangan, dan mitigasi keamanan jaringan.}{Jawaban tertulis kuis.}{Identifikasi ancaman dan serangan jaringan & 3 & Jenis ancaman (DoS, MITM, sniffing, dll.) dan mekanisme serangan dijelaskan dengan benar dan lengkap. & Sebagian besar ancaman dipahami, tetapi mekanisme atau dampak serangan masih kurang detail. & Ancaman jaringan tidak diidentifikasi dengan benar atau penjelasan tidak relevan. \\ \\
Pemahaman mekanisme perlindungan & 2 & Firewall, VPN, enkripsi, dan monitoring jaringan dijelaskan dengan benar dan kegunaannya dalam konteks perlindungan jaringan tepat. & Mekanisme perlindungan cukup dipahami, tetapi penerapannya pada konteks keamanan masih kurang spesifik. & Mekanisme perlindungan tidak dipahami atau penjelasan mengandung kesalahan mendasar. \\}

\assessmentblock{PBL Desain Jaringan}{Project Based Learning}{SCPMK0504-03201, SCPMK1002-03203}{Mahasiswa merancang infrastruktur jaringan komputer lengkap beserta analisis keamanannya untuk skenario organisasi yang diberikan.}{Menganalisis kebutuhan organisasi, merancang topologi, memilih protokol, mengkonfigurasi simulator, menganalisis keamanan, dan mempresentasikan desain.}{Dokumen desain jaringan, konfigurasi simulator, laporan keamanan, dan/atau bahan presentasi.}{Kualitas desain arsitektur jaringan & 9 & Desain mencakup topologi yang tepat, pemilihan protokol yang sesuai, pengalamatan IP yang benar, dan memenuhi kebutuhan organisasi. & Desain sebagian besar memenuhi kebutuhan, tetapi beberapa aspek protokol atau pengalamatan belum optimal. & Desain tidak memenuhi kebutuhan organisasi atau memiliki kesalahan arsitektur yang fundamental. \\ \\
Fungsionalitas konfigurasi jaringan & 7 & Seluruh konfigurasi pada simulator berjalan dengan benar, konektivitas antarsegmen terjamin, dan routing berfungsi. & Konfigurasi sebagian besar berjalan, tetapi beberapa segmen atau routing masih ada kekurangan. & Konfigurasi tidak berjalan atau banyak kesalahan yang menyebabkan konektivitas gagal. \\ \\
Komprehensivitas analisis keamanan & 6 & Analisis keamanan mencakup identifikasi ancaman, implementasi kontrol keamanan, dan rekomendasi yang spesifik dan dapat diterapkan. & Analisis keamanan cukup komprehensif, tetapi beberapa ancaman atau rekomendasi masih terlalu umum. & Analisis keamanan tidak komprehensif atau rekomendasi tidak dapat diterapkan pada desain jaringan. \\}

\assessmentblock{UAS Presentasi Desain Jaringan dan Analisis Keamanan}{UAS presentasi dan defense}{SCPMK0504-03201, SCPMK1002-03203}{Mahasiswa mempresentasikan dan mempertahankan desain jaringan dan analisis keamanan yang telah dikembangkan.}{Menyiapkan presentasi teknis, mempresentasikan desain, menjawab pertanyaan teknis, dan mengumpulkan dokumen final.}{Bahan presentasi, dokumen desain jaringan final, laporan keamanan komprehensif.}{Kualitas desain jaringan final & 10 & Desain jaringan final lengkap, semua kebutuhan organisasi terpenuhi, konfigurasi berjalan, dan dokumentasi teknis komprehensif. & Desain jaringan cukup lengkap, tetapi beberapa kebutuhan atau dokumentasi teknis masih kurang. & Desain tidak lengkap atau banyak kebutuhan organisasi yang tidak terpenuhi. \\ \\
Kedalaman analisis keamanan & 9 & Analisis keamanan mendalam, mencakup threat modeling, implementasi kontrol, dan rencana respons insiden yang realistis. & Analisis keamanan cukup mendalam, tetapi threat modeling atau rencana respons insiden masih terbatas. & Analisis keamanan dangkal atau tidak menghubungkan ancaman dengan kontrol yang diimplementasikan. \\ \\
Kemampuan presentasi dan defense teknis & 6 & Presentasi terstruktur, penjelasan teknis akurat, dan pertanyaan teknis dijawab dengan pemahaman mendalam. & Presentasi cukup terstruktur, tetapi beberapa penjelasan teknis atau jawaban pertanyaan masih kurang mendalam. & Tidak mampu menjelaskan desain secara teknis atau memberikan justifikasi keputusan desain. \\}

\begin{tblr}{width=\textwidth,colspec={|Q[l,m,3.2cm]|X[l,m]|},hlines,vlines,hspan=minimal,row{1-3}={bg=headgray},rowsep=1.5pt,colsep=3pt}
Jadwal Pelaksanaan: & Minggu 1--17 sesuai RPS. Setiap evaluasi memiliki RTM dan rubrik multi-kriteria. \\
Lain-Lain Yang Diperlukan: & LMS, Cisco Packet Tracer, Wireshark, dokumentasi jaringan, dan perangkat presentasi. \\
Pustaka: & Mengacu pustaka utama dan pendukung pada bagian RPS. \\
\end{tblr}
